% GENERATED FILE. Edit canonical lessons, then run npm run generate:books.
% canonical-source-hash: fnv1a64:4c2a63e2ed8b1fc8
% canonical-lessons: TA-C319-antu-niraivu, TA-C319-ponvila, TA-C319-antuvila, TA-C319-vilakkolam, TA-C319-alaippital

\chapter{Jubilees and Anniversaries}
\label{ch:ta-jubilees-and-anniversaries}
\hlchaptermodality{319}

\begin{chapteropening}
\textbf{By the end of this chapter:} \emph{I can talk about anniversaries and yearly celebrations.}

Complete the last lesson of chapter 319: I can talk about anniversaries and yearly celebrations.
\end{chapteropening}

\section[\texorpdfstring{\ta{ஆண்டு} \ta{நிறைவு} (āṇṭu niṟaivu)}{ஆண்டு நிறைவு (āṇṭu niṟaivu)}]{\ta{ஆண்டு} \ta{நிறைவு} (āṇṭu niṟaivu) — an anniversary}

\label{lesson:TA-C319-antu-niraivu}



\begin{quote}
Before the new one: say the Tamil for to preside, then the Tamil for to give a speech.
\end{quote}

\subsection*{You'll want to know: \ta{ஆண்டு} \ta{நிறைவு}}
\textbf{\ta{ஆண்டு} \ta{நிறைவு}} — \emph{āṇṭu niṟaivu} — \textquotedblleft{}an anniversary\textquotedblright{}.

\begin{quote}
\ta{எங்கள்} \ta{பள்ளியின்} \ta{ஐம்பதாம்} \ta{ஆண்டு} \ta{நிறைவு} \ta{அடுத்த} \ta{மாதம்}.
\end{quote}

\begin{quote}
\emph{eṅkaḷ paḷḷiyiṉ aimpatām āṇṭu niṟaivu aṭutta mātam.}
\end{quote}

\begin{quote}
Our school's fiftieth anniversary is next month.
\end{quote}

With an ordinal: \ta{பத்தாம்} \ta{ஆண்டு} \ta{நிறைவு}, the tenth anniversary.

\begin{cousinweb}
\textbf{\ta{ஆண்டு} \ta{நிறைவு}} (\emph{āṇṭu niṟaivu}) joins native \ta{ஆண்டு}, \textquotedblleft{}year\textquotedblright{} (DEDR 5153), and \ta{நிறைவு}, \textquotedblleft{}fullness, completion\textquotedblright{}, from \ta{நிறை}, \textquotedblleft{}to become full\textquotedblright{} (DEDR 3682).
\end{cousinweb}

\begin{grammarlens}[title={What you've built}]
One more word for telling what happened.
\end{grammarlens}

\subsection*{Your turn}
\begin{itemize}
\raggedright
  \item \emph{Say it:} \emph{āṇṭu niṟaivu}
  \item \emph{Say it:} \emph{āṇṭu niṟaivu}, once more
  \item \emph{Say it:} \emph{eṅkaḷ paḷḷiyiṉ aimpatām āṇṭu niṟaivu aṭutta mātam.} — Our school's fiftieth anniversary is next month.
  \item \emph{Recall:} say the Tamil for to preside, then the Tamil for to give a speech, then say \emph{āṇṭu niṟaivu} again
\end{itemize}

\begin{tcolorbox}[breakable,colback=teal!4,colframe=teal!35!black,title={Before you move on}]
What does \ta{ஆண்டு} \ta{நிறைவு} mean? (\textquotedblleft{}An anniversary\textquotedblright{}.) Say it once more.
\end{tcolorbox}
\section[\texorpdfstring{\ta{பொன்விழா} (poṉviḻā)}{பொன்விழா (poṉviḻā)}]{\ta{பொன்விழா} (poṉviḻā) — a golden jubilee}

\label{lesson:TA-C319-ponvila}



\begin{quote}
Before the new one: say the Tamil for to give a speech, then the Tamil for an anniversary.
\end{quote}

\subsection*{You'll want to know: \ta{பொன்விழா}}
\textbf{\ta{பொன்விழா}} — \emph{poṉviḻā} — \textquotedblleft{}a golden jubilee\textquotedblright{}.

\begin{quote}
\ta{கடையின்} \ta{பொன்விழாவைப்} \ta{பெரிதாகக்} \ta{கொண்டாடினார்கள்}.
\end{quote}

\begin{quote}
\emph{kaṭaiyiṉ poṉviḻāvaip peritākak koṇṭāṭiṉārkaḷ.}
\end{quote}

\begin{quote}
They celebrated the shop's golden jubilee in a big way.
\end{quote}

\ta{பொன்}, gold: fifty years; \ta{வெள்ளிவிழா}, silver, is twenty-five.

\begin{cousinweb}
\textbf{\ta{பொன்விழா}} (\emph{poṉviḻā}) is a Tamil compound built from \ta{பொன்}, \textquotedblleft{}gold\textquotedblright{} (DEDR 4570), and \ta{விழா}, \textquotedblleft{}festival\textquotedblright{} (DEDR 5427).
\end{cousinweb}

\begin{grammarlens}[title={What you've built}]
One more word for telling what happened.
\end{grammarlens}

\subsection*{Your turn}
\begin{itemize}
\raggedright
  \item \emph{Say it:} \emph{poṉviḻā}
  \item \emph{Say it:} \emph{poṉviḻā}, once more
  \item \emph{Say it:} \emph{kaṭaiyiṉ poṉviḻāvaip peritākak koṇṭāṭiṉārkaḷ.} — They celebrated the shop's golden jubilee in a big way.
  \item \emph{Recall:} say the Tamil for to give a speech, then the Tamil for an anniversary, then say \emph{poṉviḻā} again
\end{itemize}

\begin{tcolorbox}[breakable,colback=teal!4,colframe=teal!35!black,title={Before you move on}]
What does \ta{பொன்விழா} mean? (\textquotedblleft{}A golden jubilee\textquotedblright{}.) Say it once more.
\end{tcolorbox}
\section[\texorpdfstring{\ta{ஆண்டுவிழா} (āṇṭuviḻā)}{ஆண்டுவிழா (āṇṭuviḻā)}]{\ta{ஆண்டுவிழா} (āṇṭuviḻā) — an annual day}

\label{lesson:TA-C319-antuvila}



\begin{quote}
Before the new one: say the Tamil for an anniversary, then the Tamil for a golden jubilee.
\end{quote}

\subsection*{You'll want to know: \ta{ஆண்டுவிழா}}
\textbf{\ta{ஆண்டுவிழா}} — \emph{āṇṭuviḻā} — \textquotedblleft{}an annual day\textquotedblright{}.

\begin{quote}
\ta{பள்ளி} \ta{ஆண்டுவிழாவில்} \ta{என்} \ta{மகள்} \ta{நடனம்} \ta{ஆடினாள்}.
\end{quote}

\begin{quote}
\emph{paḷḷi āṇṭuviḻāvil eṉ makaḷ naṭaṉam āṭiṉāḷ.}
\end{quote}

\begin{quote}
My daughter danced at the school's annual day.
\end{quote}

A yearly celebration, often at schools and offices.

\begin{cousinweb}
\textbf{\ta{ஆண்டுவிழா}} (\emph{āṇṭuviḻā}) is a Tamil compound built from \ta{ஆண்டு}, \textquotedblleft{}year\textquotedblright{} (DEDR 5153), and \ta{விழா}, \textquotedblleft{}festival\textquotedblright{} (DEDR 5427).
\end{cousinweb}

\begin{grammarlens}[title={What you've built}]
One more word for telling what happened.
\end{grammarlens}

\subsection*{Your turn}
\begin{itemize}
\raggedright
  \item \emph{Say it:} \emph{āṇṭuviḻā}
  \item \emph{Say it:} \emph{āṇṭuviḻā}, once more
  \item \emph{Say it:} \emph{paḷḷi āṇṭuviḻāvil eṉ makaḷ naṭaṉam āṭiṉāḷ.} — My daughter danced at the school's annual day.
  \item \emph{Recall:} say the Tamil for an anniversary, then the Tamil for a golden jubilee, then say \emph{āṇṭuviḻā} again
\end{itemize}

\begin{tcolorbox}[breakable,colback=teal!4,colframe=teal!35!black,title={Before you move on}]
What does \ta{ஆண்டுவிழா} mean? (\textquotedblleft{}An annual day\textquotedblright{}.) Say it once more.
\end{tcolorbox}
\section[\texorpdfstring{\ta{விழாக்கோலம்} (viḻākkōlam)}{விழாக்கோலம் (viḻākkōlam)}]{\ta{விழாக்கோலம்} (viḻākkōlam) — a festive look}

\label{lesson:TA-C319-vilakkolam}



\begin{quote}
Before the new one: say the Tamil for a golden jubilee, then the Tamil for an annual day.
\end{quote}

\subsection*{You'll want to know: \ta{விழாக்கோலம்}}
\textbf{\ta{விழாக்கோலம்}} — \emph{viḻākkōlam} — \textquotedblleft{}a festive look\textquotedblright{}.

\begin{quote}
\ta{பண்டிகைக்கு} \ta{ஊர்} \ta{விழாக்கோலம்} \ta{பூண்டது}.
\end{quote}

\begin{quote}
\emph{paṇṭikaikku ūr viḻākkōlam pūṇṭatu.}
\end{quote}

\begin{quote}
For the festival the town put on a festive look.
\end{quote}

Usually with \ta{பூண்}, to put on: \ta{விழாக்கோலம்} \ta{பூண்டது}.

\begin{cousinweb}
\textbf{\ta{விழாக்கோலம்}} (\emph{viḻākkōlam}) is a Tamil compound built from \ta{விழா}, \textquotedblleft{}festival\textquotedblright{} (DEDR 5427), and \ta{கோலம்}, \textquotedblleft{}beauty, attire, form\textquotedblright{} (DEDR 2240).
\end{cousinweb}

\begin{grammarlens}[title={What you've built}]
One more word for telling what happened.
\end{grammarlens}

\subsection*{Your turn}
\begin{itemize}
\raggedright
  \item \emph{Say it:} \emph{viḻākkōlam}
  \item \emph{Say it:} \emph{viḻākkōlam}, once more
  \item \emph{Say it:} \emph{paṇṭikaikku ūr viḻākkōlam pūṇṭatu.} — For the festival the town put on a festive look.
  \item \emph{Recall:} say the Tamil for a golden jubilee, then the Tamil for an annual day, then say \emph{viḻākkōlam} again
\end{itemize}

\begin{tcolorbox}[breakable,colback=teal!4,colframe=teal!35!black,title={Before you move on}]
What does \ta{விழாக்கோலம்} mean? (\textquotedblleft{}A festive look\textquotedblright{}.) Say it once more.
\end{tcolorbox}
\section[\texorpdfstring{\ta{அழைப்பிதழ்} (aḻaippitaḻ)}{அழைப்பிதழ் (aḻaippitaḻ)}]{\ta{அழைப்பிதழ்} (aḻaippitaḻ) — an invitation card}

\label{lesson:TA-C319-alaippital}



\begin{quote}
Before the new one: say the Tamil for an annual day, then the Tamil for a festive look.
\end{quote}

\subsection*{You'll want to know: \ta{அழைப்பிதழ்}}
\textbf{\ta{அழைப்பிதழ்}} — \emph{aḻaippitaḻ} — \textquotedblleft{}an invitation card\textquotedblright{}.

\begin{quote}
\ta{திருமண} \ta{அழைப்பிதழ்} \ta{தபாலில்} \ta{வந்தது}.
\end{quote}

\begin{quote}
\emph{tirumaṇa aḻaippitaḻ tapālil vantatu.}
\end{quote}

\begin{quote}
The wedding invitation card came by post.
\end{quote}

\ta{அழை}, to invite + \ta{இதழ்}, leaf: a written invitation.

\begin{cousinweb}
\textbf{\ta{அழைப்பிதழ்}} (\emph{aḻaippitaḻ}) is a Tamil compound built from \ta{அழைப்பு}, \textquotedblleft{}a call, an invitation\textquotedblright{}, from \ta{அழை}, \textquotedblleft{}to call, summon\textquotedblright{} (DEDR 282), and \ta{இதழ்}, \textquotedblleft{}a leaf\textquotedblright{} (DEDR 459).
\end{cousinweb}

\begin{grammarlens}[title={What you've built}]
One more word for telling what happened.
\end{grammarlens}

\subsection*{Your turn}
\begin{itemize}
\raggedright
  \item \emph{Say it:} \emph{aḻaippitaḻ}
  \item \emph{Say it:} \emph{aḻaippitaḻ}, once more
  \item \emph{Say it:} \emph{tirumaṇa aḻaippitaḻ tapālil vantatu.} — The wedding invitation card came by post.
  \item \emph{Recall:} say the Tamil for an annual day, then the Tamil for a festive look, then say \emph{aḻaippitaḻ} again
\end{itemize}

\begin{tcolorbox}[breakable,colback=teal!4,colframe=teal!35!black,title={Before you move on}]
What does \ta{அழைப்பிதழ்} mean? (\textquotedblleft{}An invitation card\textquotedblright{}.) Say it once more.
\end{tcolorbox}
